% Part 2 of Day 13: Logging.
%
% Topics of the syllabus: structured logs, log levels, log rotation,
% sampling; time-series storage; privacy of logged data. Exercise: calculate
% the storage that one day of logs needs at three sampling rates.
%
% The numbers of this part have three origins. Each frame says so in its
% source line.
% 1. An experiment of this course: tools/experiments/day13_logging.py,
%    values in tools/data/day13/day13_logging.json. A simulated hour of the
%    Day 8 detector at 10 frames each second (synthetic records, seed 0),
%    written with the Python standard library on an x86 processor.
% 2. The Prometheus documentation (bytes for each sample, default
%    retention), copy in tools/data/day13/sources/prometheus_storage.md.
% 3. The textbook (evidence channels, sampling of successful requests and
%    errors, data minimization) and the companion book (a CSV logger, a
%    SQLite table of counts). Day 12 gave the SQLite commit times.
% The two listings ran as scripts on the work computer.

\theorypart{Logging}%
  {calculate the storage that one day of logs needs at three sampling rates}

% ------------------------------------------------------------------ channels
\begin{frame}{Four kinds of evidence}
  \footnotesize
  \renewcommand{\arraystretch}{1.15}
  \setlength{\tabcolsep}{4pt}
  \begin{tabular}{@{}L{0.18\textwidth}L{0.30\textwidth}L{0.42\textwidth}@{}}
    \toprule
    \textbf{Channel} & \textbf{Question} & \textbf{On the edge device} \\
    \midrule
    Metrics & Which value left its normal range? & Latency, frame rate,
      temperature, confidence: numbers at fixed times \\
    Logs & Which part did what, and when? & One record for each event: a
      start, an error, a slow frame \\
    Traces & Where did one input spend its time? & The time of each step
      for one frame: capture, model, send \\
    Prediction logs & Did the model behave as before? & A sample of inputs
      and outputs, for later labels \\
    \bottomrule
  \end{tabular}
  \vskip 0.6em
  \small
  \begin{itemize}
    \item Metrics start an alert. Logs and traces explain it. Prediction
          logs connect a problem to the model.
    \item On a device, all four compete for the same small disk and the
          same link. This part is about how much to keep.
  \end{itemize}
  \source{After \emph{Machine Learning Systems}, Volume II, chapter ``ML
    Operations at Scale'', section ``Observability architecture''.}
\end{frame}

% ------------------------------------------------------------------ structured
\begin{frame}[fragile]{A line for a person, a record for a program}
  \small
  Free text (110 bytes):
\begin{lstlisting}[basicstyle=\ttfamily\tiny]
2026-10-03T14:00:00.000Z INFO live_detect: frame 0, 2 objects, 47.6 ms, 9.8 fps, conf 0.608, 60.9 C, 412.1 MB
\end{lstlisting}
  Structured, one JSON object on each line (226 bytes):
\begin{lstlisting}[basicstyle=\ttfamily\tiny]
{"ts":"2026-10-03T14:00:00.000Z","level":"INFO","device":"pi-07","app":"live_detect",
 "model":"yolo11n_320_int8","event":"frame","frame":0,"latency_ms":47.6,"fps":9.8,
 "objects":2,"conf_mean":0.608,"temp_c":60.9,"rss_mb":412.1}
\end{lstlisting}
  \begin{itemize}
    \item A program reads the second form with no pattern: one call of
          \code{json.loads} for each line. A dashboard selects a field by
          its name.
    \item The free text needs a regular expression, which breaks when a
          person changes the message.
    \item The cost is the size: the field names repeat in each record.
  \end{itemize}
  \source{Experiment of this part: the first record of the simulated hour
    (the listing breaks the JSON line in three).}
\end{frame}

% ------------------------------------------------------------------ fields
\begin{frame}{The fields of a record}
  \footnotesize
  \renewcommand{\arraystretch}{1.15}
  \setlength{\tabcolsep}{4pt}
  \begin{tabular}{@{}L{0.17\textwidth}L{0.36\textwidth}L{0.37\textwidth}@{}}
    \toprule
    \textbf{Field} & \textbf{Content} & \textbf{Why} \\
    \midrule
    \code{ts} & Time in UTC, ISO 8601, with ms & Order of events, all
      devices on one time scale (Day 12) \\
    \code{level} & DEBUG, INFO, WARNING, ERROR & Filter and count \\
    \code{device}, \code{app} & Host name, program and its version & Which
      device, which release \\
    \code{model} & Model file or version & Compare two versions \\
    \code{event} & A fixed name: \code{frame}, \code{slow\_frame} & Count
      events by name \\
    Values & \code{latency\_ms}, \code{temp\_c}, \code{rss\_mb} & The unit
      is part of the name \\
    \bottomrule
  \end{tabular}
  \vskip 0.6em
  \small
  \begin{itemize}
    \item Use the same names in all programs of the system: the log, the
          MQTT message, and the dashboard.
    \item Write the time in UTC. Convert to local time only on the screen.
    \item An event name is fixed text. The values go into fields, not into
          the name.
  \end{itemize}
\end{frame}

% ------------------------------------------------------------------ levels
\begin{frame}{Log levels}
  \footnotesize
  \renewcommand{\arraystretch}{1.15}
  \setlength{\tabcolsep}{4pt}
  \begin{tabular}{@{}lrL{0.33\textwidth}L{0.33\textwidth}@{}}
    \toprule
    \textbf{Level} & \textbf{Value} & \textbf{Meaning} & \textbf{Example} \\
    \midrule
    DEBUG & 10 & Detail for a developer & The boxes of each frame \\
    INFO & 20 & Normal operation & Start, model loaded, summary \\
    WARNING & 30 & Unusual, the program continues & A frame took 300 ms \\
    ERROR & 40 & A task failed & The camera read failed \\
    CRITICAL & 50 & The program cannot continue & The model file is
      missing \\
    \bottomrule
  \end{tabular}
  \vskip 0.6em
  \small
  \begin{itemize}
    \item The program writes only the records at or above its level. A
          device runs at INFO. Set the level from a setting, so that you
          can change it with no new code.
    \item A DEBUG call at the level INFO costs 0.55 microseconds: it can
          stay in the code.
    \item A good rule: an ERROR needs action by a person; a WARNING needs a
          look when it repeats.
  \end{itemize}
  \source{Level values of the Python module \texttt{logging}. Time:
    experiment of this part, x86 processor.}
\end{frame}

% ------------------------------------------------------------------ python
\begin{frame}[fragile]{Structured logs in Python}
\begin{lstlisting}[style=python, basicstyle=\ttfamily\scriptsize]
import json, logging, logging.handlers, time

class JsonFormatter(logging.Formatter):
    def format(self, record):
        data = {"ts": time.strftime("%Y-%m-%dT%H:%M:%SZ",
                                    time.gmtime(record.created)),
                "level": record.levelname, "event": record.getMessage()}
        data.update(getattr(record, "fields", {}))
        return json.dumps(data, separators=(",", ":"))

handler = logging.handlers.RotatingFileHandler(
    "detector.jsonl", maxBytes=1024 * 1024, backupCount=5)
handler.setFormatter(JsonFormatter())
log = logging.getLogger("live_detect")
log.addHandler(handler)
log.setLevel(logging.INFO)

log.info("frame", extra={"fields": {"frame": 7, "latency_ms": 47.6}})
log.debug("boxes", extra={"fields": {"boxes": [[12, 40, 200, 300]]}})
\end{lstlisting}
  \small
  \begin{itemize}
    \item Only the standard library. The DEBUG record is not written.
    \item The handler limits the files: the next frame explains it.
  \end{itemize}
\end{frame}

% ------------------------------------------------------------------ cost
\begin{frame}{The cost of a log call}
  \begin{columns}[T]
    \begin{column}{0.50\textwidth}
      \footnotesize
      \renewcommand{\arraystretch}{1.15}
      \begin{tabular}{@{}lr@{}}
        \toprule
        \textbf{Call} & \textbf{Time} \\
        \midrule
        DEBUG at the level INFO & 0.55\,$\mu$s \\
        INFO to a JSON file & 43.3\,$\mu$s \\
        INFO through a queue & 47.8\,$\mu$s \\
        \bottomrule
      \end{tabular}
    \end{column}
    \begin{column}{0.46\textwidth}
      \small
      One record for each frame at 10 frames each second costs 0.4\,ms in
      each second: less than 0.1 percent of the time.
    \end{column}
  \end{columns}
  \vskip 0.7em
  \small
  \begin{itemize}
    \item The file handler writes the record to the disk at each call. On
          a microSD card, one write can sometimes wait much longer than the
          mean.
    \item A \code{QueueHandler} puts the record into a queue. A second
          thread (\code{QueueListener}) writes it. On one core it saves no
          time, but a slow write no longer stops the frame loop.
    \item So the time of a log call is seldom the problem. The size of the
          logs is.
  \end{itemize}
  \source{Experiment of this part: median of 5 runs of 100\,000 calls, one
    core of an x86 processor.}
\end{frame}

% ------------------------------------------------------------------ hour
\begin{frame}{One hour of logs}
  \begin{columns}[T]
    \begin{column}{0.52\textwidth}
      \footnotesize
      \renewcommand{\arraystretch}{1.15}
      \setlength{\tabcolsep}{4pt}
      \begin{tabular}{@{}lrr@{}}
        \toprule
        \textbf{Format} & \textbf{Hour} & \textbf{gzip} \\
        \midrule
        JSON lines & 7.89\,MB & 0.54\,MB \\
        CSV rows & 3.57\,MB & 0.45\,MB \\
        SQLite, row a frame & 2.03\,MB & \\
        SQLite, row a value & 6.33\,MB & \\
        \bottomrule
      \end{tabular}
    \end{column}
    \begin{column}{0.44\textwidth}
      \small
      36\,000 records: one for each frame at 10 frames each second, with
      12 errors and 180 slow frames.
    \end{column}
  \end{columns}
  \vskip 0.7em
  \small
  \begin{itemize}
    \item JSON repeats the names, so it has about twice the size of CSV. gzip
          removes most of the repetition: a factor of 14.6.
    \item A real log compresses less: real values change more than the
          simulated ones.
    \item SQLite stores numbers in binary: 59.5 bytes for a row of 7
          values. One row for each value (time, name, value) needs 30.9
          bytes for each value.
  \end{itemize}
\end{frame}

% ------------------------------------------------------------------ rotation
\begin{frame}{Rotation: a log with a size limit}
  \begingroup\centering
  \begin{tikzpicture}[font=\scriptsize, >=stealth,
      f/.style={draw=coolgray, rounded corners=2pt, fill=boxgray,
                align=center, text width=1.2cm, minimum height=0.75cm}]
    \node[f, fill=cardinalred!12, text width=1.9cm] (a) at (0,0)
      {\texttt{detector.log}\\ writing};
    \node[f] (b) at (2.35,0) {\texttt{.log.1}};
    \node[f] (c) at (4.15,0) {\texttt{.log.2}};
    \node at (5.45,0) {\dots};
    \node[f] (e) at (6.75,0) {\texttt{.log.5}};
    \node[f, fill=white, draw=coolgray!60, text=coolgray] (x) at (8.55,0)
      {deleted};
    \draw[->] (a) -- (b); \draw[->] (b) -- (c); \draw[->] (c) -- (5.15,0);
    \draw[->] (5.75,0) -- (e); \draw[->] (e) -- (x);
    \node[darkcyan] at (4.2,-0.75) {when the file reaches 1 MB, each file
      moves one place};
  \end{tikzpicture}\par\endgroup
  \vskip 0.5em
  \footnotesize
  \renewcommand{\arraystretch}{1.12}
  \begin{tabular}{@{}lrr@{}}
    \toprule
    \textbf{10 MB of records, limit 1 MB, 5 old files} & \textbf{Files} &
      \textbf{On the disk} \\
    \midrule
    Plain files & 6 & 5.0\,MB \\
    Old files with gzip & 6 & 0.31\,MB \\
    \bottomrule
  \end{tabular}
  \vskip 0.5em
  \small
  \begin{itemize}
    \item The disk use has a fixed limit, whatever the program writes. The
          oldest records go first.
    \item \code{TimedRotatingFileHandler(when="midnight", backupCount=7)}
          keeps one file for each day, for 7 days.
    \item System logs: \code{logrotate}, and \code{SystemMaxUse=} in the
          settings of \code{journald}.
  \end{itemize}
  \source{Experiment of this part. The gzip step follows the Python
    logging cookbook (\texttt{namer}, \texttt{rotator}).}
\end{frame}

% ------------------------------------------------------------------ sampling
\begin{frame}{Sampling: keep fewer records}
  \footnotesize
  \renewcommand{\arraystretch}{1.15}
  \setlength{\tabcolsep}{4pt}
  \begin{tabular}{@{}L{0.33\textwidth}rrrrr@{}}
    \toprule
    \textbf{Rule (one hour)} & \textbf{Records} & \textbf{Size} &
      \textbf{Errors} & \textbf{Slow} & \textbf{Max.\ ms} \\
    \midrule
    Each frame & 36\,000 & 7.89\,MB & 12 & 180 & 398.1 \\
    1 in 10 frames & 3600 & 807\,KB & 1 & 14 & 394.6 \\
    1 in 100 frames & 360 & 80.8\,KB & 0 & 4 & 324.5 \\
    1 in 100, each WARNING and ERROR & 548 & 124\,KB & 12 & 180 & 398.1 \\
    Summary each 60\,s, each ERROR & 72 & 18.9\,KB & 12 & 180 & 398.1 \\
    \bottomrule
  \end{tabular}
  \vskip 0.6em
  \small
  \begin{itemize}
    \item ``Slow'': the slow frames that the records show, of 180.
          ``Max.'': the largest latency that the records show.
    \item A plain sample loses the rare events: 1 in 100 kept no error and
          4 of 180 slow frames.
    \item Keep each rare event, and sample or summarize the normal ones.
          The textbook gives the same rule for a data centre: all errors,
          1 percent of the successful requests.
  \end{itemize}
  \source{Rule: \emph{Machine Learning Systems}, Volume I, chapter ``ML
    Operations''.}
\end{frame}

% ------------------------------------------------------------------ summary
\begin{frame}[fragile]{Summaries: aggregate on the device}
\begin{lstlisting}[basicstyle=\ttfamily\tiny]
{"ts":"2026-10-03T14:01:00.000Z","level":"INFO","device":"pi-07","app":"live_detect",
 "model":"yolo11n_320_int8","event":"summary_60s","frames":600,"latency_mean_ms":47.1,
 "latency_p95_ms":52.0,"latency_max_ms":198.4,"slow":1,"objects_mean":1.88,
 "conf_mean":0.697,"temp_c":60.6,"rss_mb":411.9}
\end{lstlisting}
  \small
  \begin{itemize}
    \item One record of 292 bytes for 600 frames. It keeps what a person
          looks at: the mean, the p95, the maximum, and the count of slow
          frames.
    \item It loses the single frame. A slow frame or an error keeps its own
          record, so you can find its time.
    \item A repeated error (a camera that fails 10 times each second) gives
          one record and then a count each minute: ``failed 600 times in
          60 s''.
    \item The lab sends the same summary with MQTT to the dashboard.
  \end{itemize}
  \source{Experiment of this part: the first summary of the simulated hour.}
\end{frame}

% ------------------------------------------------------------------ series
\begin{frame}[fragile]{Metrics are time series}
\begin{lstlisting}[basicstyle=\ttfamily\scriptsize]
edgeai_latency_ms{device="pi-07",stage="inference"}  47.6   at 14:00:00
edgeai_latency_ms{device="pi-07",stage="inference"}  46.9   at 14:00:15
edgeai_cpu_temp_celsius{device="pi-07"}              60.9   at 14:00:00
\end{lstlisting}
  \small
  \begin{itemize}
    \item A sample is a time, a name, some labels, and one number. A name
          with one set of labels is one \textbf{series}.
    \item The labels say where the value comes from: the device, the step,
          the model. A query selects or groups by label.
    \item Each new label value makes a new series. Do not put a frame
          number, an image name, or a person in a label: millions of series
          fill the storage.
    \item A log record can hold any field. A metric holds only numbers with
          a fixed set of labels, so it is small and fast to query.
  \end{itemize}
  \source{Notation of Prometheus; metric names of HW-08.}
\end{frame}

% ------------------------------------------------------------------ storage
\begin{frame}{Three ways to store metrics on a device}
  \footnotesize
  \renewcommand{\arraystretch}{1.15}
  \setlength{\tabcolsep}{4pt}
  \begin{tabular}{@{}L{0.20\textwidth}L{0.22\textwidth}L{0.24\textwidth}L{0.24\textwidth}@{}}
    \toprule
    \textbf{Store} & \textbf{Size of one frame} & \textbf{Good for} &
      \textbf{Limit} \\
    \midrule
    CSV file each day & 104.7 bytes & A simple logger, a spreadsheet & A
      query reads the whole file \\
    SQLite table & 59.5 bytes & Queries by time, one file & One writer;
      commit in batches \\
    Time-series database & 1 to 2 bytes for each value & Many series,
      retention, a query language & One more service to run \\
    \bottomrule
  \end{tabular}
  \vskip 0.6em
  \small
  \begin{itemize}
    \item The companion book logs sensor readings to a CSV file each 60 s,
          and the bee counter writes one row each 10 s to SQLite.
    \item Prometheus is a time-series database: it compresses the samples
          and deletes samples older than its retention time (default 15
          days).
    \item The lab uses the Python dashboard of HW-08 with CSV files. Option A
          of HW-08 uses Prometheus.
  \end{itemize}
  \source{Experiment of this part (CSV: the 13 fields of the record;
    SQLite: 7 numbers).\\ Prometheus documentation, ``Storage''.}
\end{frame}

% ------------------------------------------------------------------ sqlite
\begin{frame}[fragile]{SQLite for metrics}
\begin{lstlisting}[style=python, basicstyle=\ttfamily\scriptsize]
import sqlite3, time

con = sqlite3.connect("metrics.db")
con.execute("PRAGMA journal_mode=WAL")
con.execute("PRAGMA synchronous=NORMAL")
con.execute("""CREATE TABLE IF NOT EXISTS frames (t REAL, latency_ms REAL,
               fps REAL, objects INTEGER, conf_mean REAL, temp_c REAL)""")
con.execute("CREATE INDEX IF NOT EXISTS frames_t ON frames (t)")

batch = []                           # rows of the last second
batch.append((time.time(), 47.6, 9.8, 2, 0.608, 60.9))
con.executemany("INSERT INTO frames VALUES (?,?,?,?,?,?)", batch)
con.commit()                         # one commit for each second, not each frame
rows = con.execute("SELECT avg(latency_ms), max(latency_ms) FROM frames "
                   "WHERE t > ?", (time.time() - 3600,)).fetchall()
\end{lstlisting}
  \small
  \begin{itemize}
    \item Day 12 measured one commit: 2.8 ms with the defaults, 0.027 ms
          with WAL and \code{synchronous=NORMAL}.
    \item The index on the time makes a query of the last hour fast.
          Delete old rows each night: \code{DELETE ... WHERE t < ?}.
  \end{itemize}
\end{frame}

% ------------------------------------------------------------------ tiers
\begin{frame}{Keep detail for a short time}
  \begin{columns}[T]
    \begin{column}{0.56\textwidth}
      \footnotesize
      \renewcommand{\arraystretch}{1.15}
      \setlength{\tabcolsep}{4pt}
      \begin{tabular}{@{}lrr@{}}
        \toprule
        \textbf{Tier} & \textbf{Rows} & \textbf{Size} \\
        \midrule
        Each frame, 1 day & 864\,000 & 49.0\,MB \\
        Each second, 30 days & 2\,592\,000 & 147.0\,MB \\
        Each minute, 365 days & 525\,600 & 29.8\,MB \\
        \midrule
        Total & & 225.9\,MB \\
        Each frame, 365 days & & 17\,889\,MB \\
        \bottomrule
      \end{tabular}
    \end{column}
    \begin{column}{0.40\textwidth}
      \small
      A row of the SQLite table: 59.5 bytes. A summary row has about the
      same size.
    \end{column}
  \end{columns}
  \vskip 0.7em
  \small
  \begin{itemize}
    \item Recent data needs detail: you look for one slow frame of today.
          Old data needs only trends: the mean latency of last March.
    \item A job makes the summaries (downsampling) and deletes the detail
          when it is older than its tier.
    \item Prometheus: 20 series, each 15 s, for 15 days, is 1\,728\,000
          samples: 1.6 to 3.3 MB at 1 to 2 bytes each.
  \end{itemize}
  \source{Calculation with the row size of the experiment; Prometheus
    formula of its documentation. 1 MB = 1024 KB.}
\end{frame}

% ------------------------------------------------------------------ sdcard
\begin{frame}{The microSD card}
  \small
  \begin{itemize}
    \item The Raspberry Pi stores its system and its logs on a microSD card.
          A full card stops the logs, the message queue of Day 12, and
          sometimes the system.
    \item Flash memory wears with each write. Many small writes (one write
          with a flush for each frame) wear it more than one large write
          each second.
    \item The rules of this part protect the card: a size limit
          (rotation), fewer records (sampling, summaries), and writes in
          batches.
    \item Check the space: \code{df -h /}, \code{du -sh logs},
          \code{journalctl --disk-usage}.
    \item An option: write the detailed log to a RAM disk (\code{/run} or a
          \code{tmpfs}) and keep only the summaries on the card. A restart
          deletes the detail.
  \end{itemize}
  \keypoint{Give each log a limit in bytes. A log with no limit fills the
    disk on a day when nobody watches.}
\end{frame}

% ------------------------------------------------------------------ privacy
\begin{frame}{Privacy of logged data}
  \footnotesize
  \renewcommand{\arraystretch}{1.15}
  \setlength{\tabcolsep}{4pt}
  \begin{tabular}{@{}L{0.30\textwidth}L{0.30\textwidth}L{0.30\textwidth}@{}}
    \toprule
    \textbf{Do not log} & \textbf{Risk} & \textbf{Log in its place} \\
    \midrule
    Images, audio & Faces, voices, homes & Counts, classes, scores \\
    Text of a user (Day 10) & Names, health, opinions & Length, intent
      class \\
    Passwords, keys & Access to the system & Nothing; a key ID \\
    Exact place with time & The habits of a person & A zone, a coarse
      time \\
    \bottomrule
  \end{tabular}
  \vskip 0.6em
  \small
  \begin{itemize}
    \item \textbf{Data minimization}: log only what the purpose needs. A
          person counter needs a count, not the image of each person.
    \item \textbf{Purpose and retention}: write why you keep each field and
          for how long. Rotation enforces the time.
    \item \textbf{Access}: a log file on the device is readable by its
          users. Set the file rights; do not send logs over an open link.
    \item An edge device can keep the raw data local. A log that copies it
          to a server removes this advantage.
  \end{itemize}
  \source{\emph{Machine Learning Systems}, Volume I, chapter ``Responsible
    Engineering''; Volume II, chapter ``Edge Intelligence'' (GDPR).}
\end{frame}

% ------------------------------------------------------------------ prediction
\begin{frame}{Prediction logs: a sample with care}
  \small
  \begin{itemize}
    \item Part 3 needs examples of the inputs when the model is not sure:
          they become the new training data.
    \item Keep a small sample: for example, the frames with a confidence
          below a threshold, at most N each hour.
    \item Store them in a separate folder with its own size limit and a
          short retention. Do not mix them with the normal log.
    \item Images of people need a notice in the room, or a blur of the faces
          on the device before the image is stored.
    \item Record with each sample: the time, the model version, the output,
          and the reason for the sample.
  \end{itemize}
  \vskip 0.4em
  \begingroup\centering
  \begin{tikzpicture}[font=\scriptsize, >=stealth,
      st/.style={draw=coolgray, rounded corners=2pt, fill=boxgray,
                 align=center, text width=1.9cm, minimum height=0.8cm}]
    \node[st] (a) at (0,0) {Frame and output};
    \node[st, fill=cardinalred!12] (b) at (2.7,0) {Confidence\\ below 0.4?};
    \node[st] (c) at (5.4,0) {Hourly quota\\ not full?};
    \node[st, fill=darkcyan!15] (d) at (8.1,0) {Save image\\ and record};
    \draw[->] (a) -- (b); \draw[->] (b) -- node[above] {yes} (c);
    \draw[->] (c) -- node[above] {yes} (d);
  \end{tikzpicture}\par\endgroup
\end{frame}

% ------------------------------------------------------------------ lab
\begin{frame}{The logs of the lab}
  \begingroup\centering
  \begin{tikzpicture}[font=\scriptsize, >=stealth,
      st/.style={draw=coolgray, rounded corners=2pt, fill=boxgray,
                 align=center, text width=2.0cm, minimum height=0.9cm}]
    \node[st, fill=cardinalred!12] (d) at (0,0) {Day 8 detector\\ with
      measures};
    \node[st] (l) at (3.4,0.8) {JSON log\\ 1 MB x 5 files};
    \node[st] (m) at (3.4,-0.8) {Summary each 10 s\\ MQTT};
    \node[st, fill=darkcyan!15] (s) at (6.8,-0.8) {Dashboard:\\ CSV each
      day};
    \draw[->] (d) -- node[above, sloped] {events} (l);
    \draw[->] (d) -- node[below, sloped] {metrics} (m);
    \draw[->] (m) -- (s);
  \end{tikzpicture}\par\endgroup
  \vskip 0.6em
  \small
  \begin{itemize}
    \item Part A: the detector writes a structured JSON log with rotation:
          the start, each WARNING and ERROR, and a summary.
    \item Part B: the same summary goes as an MQTT message to the topic
          \code{edgeai/gNN/pi/metrics}. The dashboard of HW-08 stores it in
          a CSV file for each day and draws it.
    \item The lab keeps no image. The drift check of Part C uses the
          confidence and the input statistics.
  \end{itemize}
\end{frame}

% ------------------------------------------------------------------ exercise
\exerciseframe{Storage of one day of logs}{%
  The detector of the lab writes JSON records of 230 bytes. Calculate the
  storage for one day (1 MB = 1024 KB) for each rate:
  \begin{enumerate}
    \item[A.] one record for each frame, at 10 frames each second;
    \item[B.] one record each second;
    \item[C.] one summary each minute.
  \end{enumerate}
  The device has a log budget of 1 GB for 30 days. Which rates fit? Which
  fit when gzip compresses the old files by a factor of 14.6?}

\answerframe{Storage of one day of logs}{%
  \footnotesize
  \setlength{\tabcolsep}{4pt}
  \begin{tabular}{@{}lrrrr@{}}
    \toprule
    \textbf{Rate} & \textbf{Records} & \textbf{One day} &
      \textbf{30 days} & \textbf{30 days, gzip} \\
    \midrule
    A. Each frame & 864\,000 & 189.5\,MB & 5685\,MB & 389.4\,MB \\
    B. Each second & 86\,400 & 19.0\,MB & 568.5\,MB & 38.9\,MB \\
    C. Each minute & 1440 & 0.32\,MB & 9.5\,MB & 0.65\,MB \\
    \bottomrule
  \end{tabular}
  \vskip 0.6em
  \small
  \begin{itemize}
    \item One day: $864\,000 \times 230 = 198\,720\,000$ bytes $=$ 189.5 MB.
    \item Without gzip, B and C fit 1 GB (1024 MB) for 30 days. With gzip,
          A also fits.
    \item The factor 14.6 came from synthetic records. Measure the factor
          on real logs before you plan with it.
  \end{itemize}}
